\heading{经济学原理（II）-期末样题}
\instructions{满分 100 分，建议答题时间 120 分钟。半开卷；计算题请写出关键中间步骤，并清晰标示最终答案。}
\newcommand{\finalsamplefigdir}{figures}
\newcommand{\acctcell}[1]{\begin{tabular}[t]{@{}l@{}}#1\end{tabular}}
\newcommand{\emptyacct}[1]{%
\begin{center}
\begin{tabular}{|p{0.45\linewidth}|p{0.45\linewidth}|}
\hline
\multicolumn{2}{|c|}{#1}\\
\hline
资产 & 负债\\
\hline
\vspace{1.1cm} & \vspace{1.1cm}\\
\hline
\end{tabular}
\end{center}}

\textbf{一、选择题（16 分，每小题 2 分）}

\begin{question}
下列哪项符合丧志工人的定义？
\begin{choiceoptions}[1]
  \optionitem{过去 1 年找工作但最近 4 周没有找工作的人}
  \optionitem{想找工作但放弃找工作的人}
  \optionitem{想全职工作但目前在兼职工作的人}
  \optionitem{想上班摸鱼但害怕被领导开除的人}
\end{choiceoptions}
\begin{solution}
丧志工人是想工作、也愿意工作，但因为认为找不到合适工作而停止主动寻找工作的人。\anstext{B}
\end{solution}
\end{question}

\begin{question}
劳动参与率的计算公式是什么？
\begin{choiceoptions}[1]
  \optionitem{劳动力人口/总人口 $\times 100\%$}
  \optionitem{就业人口/成人人口 $\times 100\%$}
  \optionitem{就业人口/总人口 $\times 100\%$}
  \optionitem{劳动力人口/满 16 岁的人口 $\times 100\%$}
\end{choiceoptions}
\begin{solution}
劳动参与率等于劳动力人口占成年人口的比例；题中以满 16 岁人口作为成年人口口径。\anstext{D}
\end{solution}
\end{question}

\begin{question}
以下哪项不一定是非周期性失业？
\begin{choiceoptions}[2]
  \optionitem{结构性失业}
  \optionitem{自然失业}
  \optionitem{自愿失业}
  \optionitem{摩擦性失业}
\end{choiceoptions}
\begin{solution}
摩擦性失业和结构性失业属于非周期性失业，自然失业率也由非周期性失业构成；“自愿失业”并不必然对应非周期性失业的标准分类。\anstext{C}
\end{solution}
\end{question}

\begin{question}
以下哪一项包括在中国的 M2 但不包括在美国的 M2 中？
\begin{choiceoptions}[2]
  \optionitem{非零售货币基金}
  \optionitem{小额定期存款}
  \optionitem{活期存款}
  \optionitem{长期债券}
\end{choiceoptions}
\begin{solution}
按课堂比较口径，美国 M2 不包括非零售货币基金，而中国广义货币统计中相关口径更宽。小额定期存款和活期存款通常都在 M2 口径内，长期债券则不属于货币。\anstext{A}
\end{solution}
\end{question}

\begin{question}
美联储和中国人民银行的领导班子都是 7 人，从人数配置上来看，美联储“委员”相当于中国人民银行的什么职位？
\begin{choiceoptions}[2]
  \optionitem{行长}
  \optionitem{副行长}
  \optionitem{党委委员}
  \optionitem{党委书记}
\end{choiceoptions}
\begin{solution}
美联储理事会有 7 人：1 位主席、1 位副主席、5 位委员；中国人民银行领导班子也 7 人：1 位行长、1 位纪检监察组组长、5 位副行长，并称这种“1+1+5”的结构与美联储类似。所以从人数配置看，美联储的 5 位“委员”对应中国人民银行的 5 位副行长。\anstext{B}
\end{solution}
\end{question}

\begin{question}
Robert Lucas 因在什么方面的贡献而获得了 1995 年的诺贝尔经济学奖？
\begin{choiceoptions}[1]
  \optionitem{对货币主义的批判}
  \optionitem{分析了宏观政策中的跨期选择问题}
  \optionitem{加深了我们对金融泡沫的理解}
  \optionitem{发展和应用了理性预期假说}
\end{choiceoptions}
\begin{solution}
Lucas 的核心贡献是发展并应用理性预期假说，改变了宏观经济政策分析方式。\anstext{D}
\end{solution}
\end{question}

\begin{question}
一瓶 250ml 的椰树牌椰汁在 A 国卖 5 阿币，在 B 国卖 10 比币。根据 PPP 理论，若不考虑关税，且站在 B 国的视角分析，第一种名义汇率在长期应该是多少？
\begin{choiceoptions}[2]
  \optionitem{5}
  \optionitem{0.5}
  \optionitem{2}
  \optionitem{10}
\end{choiceoptions}
\begin{solution}
第一种名义汇率表示“1 单位外币可兑换多少单位本币”。站在 B 国视角，本币是比币，外币是阿币。根据绝对购买力平价，同一瓶椰汁换成本币后的价格应相等：
\[
 e=\frac{P_B}{P_A}=\frac{10\text{ 比币}}{5\text{ 阿币}}=2\text{ 比币/阿币}.
\]
\anstext{C}
\end{solution}
\end{question}

\begin{question}
量化宽松主要是通过以下哪项实现的？
\begin{choiceoptions}[1]
  \optionitem{央行在开放市场上大量购买国库券}
  \optionitem{央行在开放市场上大量购买金融资产}
  \optionitem{央行大量印钞}
  \optionitem{央行对利率进行严格的限制}
\end{choiceoptions}
\begin{solution}
量化宽松是央行大规模购买金融资产来扩张资产负债表、压低长期利率和改善流动性，购买对象不必限于短期国库券。\anstext{B}
\end{solution}
\end{question}

\textbf{二、中国不同代际流动人口劳动力市场现状（10 分）}

\begin{question}
改革开放以来，流动人口对中国经济发展有巨大贡献。然而，在“第四次工业革命”的冲击下，流动人口群体的就业状况令人担忧。一项基于 2016 年中国流动人口动态监测调查数据的研究发现，第二代流动人口（即父母有过外出务工或经商经历的人群）相比同时期第一代的劳动参与率较低。请根据图 1，回答问题：

\begin{center}
\includegraphics[width=0.48\linewidth]{\finalsamplefigdir/经济学原理（II）-期末样题-fig01.png}

\textbf{图 1：不同代际流动人口的劳动参与率}
\end{center}

\begin{enumerate}[label=(\alph*),leftmargin=2.5em]
  \item 整体来看，第二代流动人口的劳动参与率显著低于第一代的情况主要体现在哪个年龄段？（2 分）
  \item 图 1 显示，在某个年龄段，第二代流动人口的劳动参与率略高于第一代。额外的数据分析发现，第二代流动人口的平均受教育水平较高。请试图从机会成本的角度解释这一现象。（3 分）
  \item 在（b）问中，用机会成本解释更高的劳动参与率，与（a）问中的现象有矛盾吗？为什么？（3 分）
  \item 若丧志工人和准待业工人增加了，在第一代和第二代成人流动人口数不变的情况下，可能会对劳动参与率造成什么影响？（2 分）
\end{enumerate}
\begin{solution}
\begin{enumerate}[label=(\alph*),leftmargin=2.5em]
  \item 主要体现在较年轻年龄段，尤其是 16--25 岁附近；其中 16--20 岁左右的差距最明显。
  \item 第二代流动人口平均教育水平较高，完成教育并进入劳动力市场后，其潜在工资通常更高。工资越高，不参加工作的机会成本越高，因此在一些成年年龄段，第二代更有动力进入劳动力市场，劳动参与率可能略高。
  \item 不矛盾。年轻阶段，受教育水平较高可能意味着继续上学或接受培训的收益更大，因而工作的机会成本较高，劳动参与率偏低；完成教育后，较高人力资本又会提高潜在工资，使“不工作”的机会成本上升，从而提高劳动参与率。同一个教育因素在不同生命周期阶段可能产生不同方向的影响。
  \item 丧志工人和准待业工人通常不被计入劳动力人口。若成人流动人口数不变而这类人增加，劳动力人口占比下降，劳动参与率会下降。
\end{enumerate}
\end{solution}
\end{question}

\textbf{三、21 世纪美国失业率的变化趋势（7 分）}

\begin{question}
图 2 展示了来自美国劳工统计局的实际失业率和来自美国国会预算办公室的非周期性失业率 2000 年第一季度至 2026 年第一季度的数据。请根据图 2，回答如下问题：

\begin{center}
\includegraphics[width=0.94\linewidth]{\finalsamplefigdir/经济学原理（II）-期末样题-fig02.png}

\textbf{图 2：美国 2000Q1--2026Q1 的实际失业率（实线）与非周期性失业率（虚线）}
\end{center}

\begin{enumerate}[label=(\alph*),leftmargin=2.5em]
  \item 请描述在此期间内周期性失业率的变化趋势，注意结合经济周期和周期性失业率的符号讨论。（3 分）
  \item 请提供一个例子解释为什么有一些结构性失业对经济是好的事情。（2 分）
  \item 提供一种减少非周期性失业率的方法或例子，以及一个减少周期性失业率的办法或政策。（2 分）
\end{enumerate}
\begin{solution}
\begin{enumerate}[label=(\alph*),leftmargin=2.5em]
  \item 周期性失业率应按课堂定义理解为实际失业率相对于非周期性（自然）失业率的偏离：
  \[
  u^{cyc}=u-u^{noncyc}.
  \]
  因此，实线高于虚线时，周期性失业率为正，表示经济处在需求不足、衰退或复苏早期；实线低于虚线时，周期性失业率为负，表示经济高于正常就业水平、劳动力市场偏热。图中 2000 年前后实际失业率低于虚线，周期性失业率为负；2001 年衰退后转正并在 2002--2004 年为正；2006--2007 年又转为负。2008--2009 年金融危机后周期性失业率大幅转正，随后多年逐步回落，并在 2017--2019 年转为负。2020 年疫情冲击使周期性失业率短暂大幅为正，之后迅速下降，2022 年后多为负或接近零，2025--2026 年附近大体回到零附近。
  \item 例如技术进步淘汰低效率岗位，使劳动者从低生产率行业转向新能源、数字经济、高端制造等更高生产率行业。短期内这会造成结构性失业，但长期有助于资源重新配置、提高社会总产出。
  \item 减少非周期性失业率的办法包括职业培训、就业信息平台、降低跨地区流动成本、改善岗位匹配等。减少周期性失业率的办法包括扩张性财政政策、扩张性货币政策，例如增加政府支出、减税、降息或公开市场买入债券以刺激总需求。
\end{enumerate}
\end{solution}
\end{question}

\textbf{四、个人交易行为与银行系统（14 分）}

\begin{question}
假设小吴将藏在枕头底下的 6 万元现金存入了自己的工商银行账户中，再假设央行设置的法定准备金率为 5\%，请回答如下问题：
\begin{enumerate}[label=(\alph*),leftmargin=2.5em]
  \item 请展示小吴的存款行为对工商银行 T 账户的初始影响（即工商银行还没有对小吴的行为做出下一步操作时的影响）。请区分法定准备金和超额准备金。（2 分）
\end{enumerate}
\emptyacct{工商银行的资产负债表}
\begin{enumerate}[label=(\alph*),leftmargin=2.5em,resume]
  \item 在（a）问中，货币供给的变化是多少？基础货币的变化是多少？（2 分）
  \item 假设工商银行在合法范围内最大化贷款，并且这些贷款暂时不会回到自己手上，请展示工商银行 T 账户目前的变化情况。（2 分）
\end{enumerate}
\emptyacct{工商银行的资产负债表}
\begin{enumerate}[label=(\alph*),leftmargin=2.5em,resume]
  \item 在（c）问中，货币供给的变化是多少？基础货币的变化是多少？（2 分）
  \item 如果新的贷款最终都会全数存入银行系统，而新的超额准备金又会被全数贷出，则小吴的行为对银行系统（包括工商银行）的总连锁效应是什么？请用 T 账户来展示效应。（3 分）
\end{enumerate}
\emptyacct{银行系统的资产负债表}
\begin{enumerate}[label=(\alph*),leftmargin=2.5em,resume]
  \item 最后，假设央行卖给小吴 5000 元国库券，并从其工商银行账户中扣款。请展示该交易对工商银行 T 账户的初始影响（假设在银行系统的连锁反应中，工商银行都没有被牵涉）。按照规定，工商银行接下来要怎么做？（3 分）
\end{enumerate}
\emptyacct{工商银行的资产负债表}
\begin{solution}
\begin{enumerate}[label=(\alph*),leftmargin=2.5em]
  \item 存款 60,000 元，法定准备金率为 5\%，所以法定准备金为 $60,000\times 5\%=3,000$ 元，超额准备金为 $57,000$ 元。
\end{enumerate}
\begin{center}
\begin{tabular}{|p{0.45\linewidth}|p{0.45\linewidth}|}
\hline
\multicolumn{2}{|c|}{工商银行的资产负债表（变化量）}\\
\hline
资产 & 负债\\
\hline
\acctcell{准备金 $+60,000$\\\quad 法定准备金 $+3,000$\\\quad 超额准备金 $+57,000$} & \acctcell{存款 $+60,000$}\\
\hline
\end{tabular}
\end{center}
\begin{enumerate}[label=(\alph*),leftmargin=2.5em,start=2]
  \item 现金从公众手中转为银行存款，货币供给中“现金”减少 60,000 而“存款”增加 60,000，因此货币供给不变。基础货币中“流通中现金”减少 60,000 而“银行准备金”增加 60,000，因此基础货币也不变。
  \item 工商银行最多贷出超额准备金 57,000 元，贷款暂时不回到工商银行，因此该行保留 3,000 元准备金以满足原存款的法定准备金要求。
\end{enumerate}
\begin{center}
\begin{tabular}{|p{0.45\linewidth}|p{0.45\linewidth}|}
\hline
\multicolumn{2}{|c|}{工商银行的资产负债表（变化量）}\\
\hline
资产 & 负债\\
\hline
\acctcell{准备金 $+3,000$\\贷款 $+57,000$} & \acctcell{存款 $+60,000$}\\
\hline
\end{tabular}
\end{center}
\begin{enumerate}[label=(\alph*),leftmargin=2.5em,start=4]
  \item 与小吴存款前相比，货币供给增加 57,000 元，因为新增贷款最终形成公众可使用的新增存款或现金；基础货币仍不变，因为商业银行贷款不创造基础货币。
  \item 银行系统的存款乘数为
  \[
  \frac{1}{rr}=\frac{1}{0.05}=20.
  \]
  60,000 元新增准备金最终支持总存款 $60,000/0.05=1,200,000$ 元，其中总贷款为 $1,200,000-60,000=1,140,000$ 元。
\end{enumerate}
\begin{center}
\begin{tabular}{|p{0.45\linewidth}|p{0.45\linewidth}|}
\hline
\multicolumn{2}{|c|}{银行系统的资产负债表（总变化量）}\\
\hline
资产 & 负债\\
\hline
\acctcell{准备金 $+60,000$\\贷款 $+1,140,000$} & \acctcell{存款 $+1,200,000$}\\
\hline
\end{tabular}
\end{center}
\begin{enumerate}[label=(\alph*),leftmargin=2.5em,start=6]
  \item 央行从小吴工商银行账户扣款会使工商银行存款负债减少 5,000 元，同时准备金被央行划走，资产端准备金减少 5,000 元。
\end{enumerate}
\begin{center}
\begin{tabular}{|p{0.45\linewidth}|p{0.45\linewidth}|}
\hline
\multicolumn{2}{|c|}{工商银行的资产负债表（初始变化量）}\\
\hline
资产 & 负债\\
\hline
\acctcell{准备金 $-5,000$} & \acctcell{存款 $-5,000$}\\
\hline
\end{tabular}
\end{center}
扣款后工商银行存款余额为 55,000 元，法定准备金要求为 $55,000\times 5\%=2,750$ 元。若它此前只有 3,000 元准备金，支付 5,000 元后将出现准备金缺口；相对于合规要求，缺口为 $2,750-(-2,000)=4,750$ 元。因此它需要通过借入准备金、出售资产、收回或减少贷款等方式补足准备金。
\end{solution}
\end{question}

\textbf{五、货币数量论与通货膨胀（10 分）}

\begin{question}
如果真实货币需求函数 $(M/P)^d=L(i,Y)=Y/(4i)$，请回答如下问题：
\begin{enumerate}[label=(\alph*),leftmargin=2.5em]
  \item 如果产出增长率为 $g$，名义利率不变，请计算真实货币需求的增长率。（2 分）
  \item 货币流通速度是多少？（2 分）
  \item 如果通货膨胀率和名义利率不变，货币流通速度的增长率是多少？（2 分）
  \item 永久性（一劳永逸）地提高名义利率会如何影响货币流通速度？货币流通速度在其后的增长率又会如何受到影响？（2 分）
  \item 如果央行希望盯住长期通胀目标 $\pi$，货币供给的增速应该设为多少？（2 分）
\end{enumerate}
\begin{solution}
\begin{enumerate}[label=(\alph*),leftmargin=2.5em]
  \item 当 $i$ 不变时，$(M/P)^d=Y/(4i)$ 与 $Y$ 等比例变化，因此真实货币需求增长率为 $g$。
  \item 货币市场出清时 $M/P=Y/(4i)$，所以
  \[
  M=\frac{PY}{4i},\qquad V=\frac{PY}{M}=4i.
  \]
  \item 若 $i$ 不变，则 $V=4i$ 不变，货币流通速度增长率为 $0$。
  \item 一次性永久提高 $i$ 会使 $V=4i$ 的水平一次性上升；之后若名义利率维持在新的常数水平，$V$ 的增长率仍为 $0$。
  \item 长期中 $V$ 不变时，数量方程 $MV=PY$ 给出
  \[
  \mu+0=\pi+g,
  \]
  因此货币供给增速应为
  \[
  \mu=\pi+g.
  \]
\end{enumerate}
\ansmath{V=4i,\quad \mu=\pi+g}
\end{solution}
\end{question}

\textbf{六、国际收支平衡表与国际贸易（14 分）}

\begin{question}
请将下表补充完整（6 分），并回答表格下面的问题：

\begin{center}
\small
\begin{tabular}{lrrr}
\hline
项目 & 差额 & 贷方/负债 & 借方/资产\\
\hline
1. 经常项目 & \blank[1.5cm] & 310,000,000 & -260,000,000\\
1.A. 货物和服务 & \blank[1.5cm] & 290,000,000 & -230,000,000\\
1.A.a. 货物 & 50,000,000 & 260,000,000 & \blank[1.5cm]\\
1.A.b. 服务 & \blank[1.5cm] & 30,000,000 & \blank[1.5cm]\\
1.B. 初次收入 & -11,000,000 & \blank[1.5cm] & -27,000,000\\
1.C. 二次转移 & 1,000,000 & 4,000,000 & \blank[1.5cm]\\
2. 资本和金融账户 & -20,000,000 & -- & --\\
2.1. 资本账户 & 10,000 & 40,000 & \blank[1.5cm]\\
2.2. 金融账户 & \blank[1.5cm] & \blank[1.5cm] & -90,000,000\\
2.2.1. 非储备性质 & 3,000,000 & 69,990,000 & \blank[1.5cm]\\
2.2.2. 储备资产 & -23,010,000 & -- & --\\
3. 净误差与遗漏 & \blank[1.5cm] & -- & --\\
\hline
\end{tabular}
\end{center}

\begin{enumerate}[label=(\alph*),leftmargin=2.5em]
  \item 该国净误差与遗漏绝对值占进、出口总额的百分比大约是多少？这说明该国数据失真的严重程度高吗？（2 分）
  \item 如果该国净资本流入（NCI）统计不存在任何误差与遗漏，则理论上净出口（NX）应该等于多少？（2 分）
  \item 该国的货币被称为 X 币，且在外汇市场上 1 美元可以兑换 2 单位 X 币，假如 X 国一件标价为 690 单位 X 币的金属饰品在美国标价为 390 美元，那么以这件金属饰品来衡量的该国相对于美国产品的真实汇率是多少？假设该产品在两国之间运输的成本可忽略不计，且两国对此产品不征税，请使用课上讲的第一种汇率的表示方式。（2 分）
  \item 假如该国政府通过外汇管制提高名义汇率，使得 1 美元可以兑换 3 单位 X 币，但与此同时该国发生了通货膨胀，导致该金属制品在 X 国标价为 790 单位 X 币，而美国市场上该产品的价格则下降到了 310 美元。请问，X 国生产的这种产品的国际竞争力是上升了还是下降了？为什么？（2 分）
\end{enumerate}
\begin{solution}
表格补全如下。差额按“贷方/负债减借方/资产”的符号口径填写。
\begin{center}
\small
\begin{tabular}{lrrr}
\hline
项目 & 差额 & 贷方/负债 & 借方/资产\\
\hline
1. 经常项目 & 50,000,000 & 310,000,000 & -260,000,000\\
1.A. 货物和服务 & 60,000,000 & 290,000,000 & -230,000,000\\
1.A.a. 货物 & 50,000,000 & 260,000,000 & -210,000,000\\
1.A.b. 服务 & 10,000,000 & 30,000,000 & -20,000,000\\
1.B. 初次收入 & -11,000,000 & 16,000,000 & -27,000,000\\
1.C. 二次转移 & 1,000,000 & 4,000,000 & -3,000,000\\
2. 资本和金融账户 & -20,000,000 & -- & --\\
2.1. 资本账户 & 10,000 & 40,000 & -30,000\\
2.2. 金融账户 & -20,010,000 & 69,990,000 & -90,000,000\\
2.2.1. 非储备性质 & 3,000,000 & 69,990,000 & -66,990,000\\
2.2.2. 储备资产 & -23,010,000 & -- & --\\
3. 净误差与遗漏 & -30,000,000 & -- & --\\
\hline
\end{tabular}
\end{center}
\begin{enumerate}[label=(\alph*),leftmargin=2.5em]
  \item 经常项目差额为 50,000,000，资本和金融账户差额为 $-20,000,000$，因此净误差与遗漏为
  \[
  -\bigl(50,000,000-20,000,000\bigr)=-30,000,000.
  \]
  按讲义中对该类表格采用的广义进、出口总额口径，分母为经常项目和资本账户的贷方、借方绝对值之和：
  \[
  310,000,000+260,000,000+40,000+30,000=570,070,000.
  \]
  所以
  \[
  \frac{|{-30,000,000}|}{570,070,000}\approx 5.26\%.
  \]
  这个比例明显不低，说明数据失真或统计误差比较值得重视。若只用货物和服务进出口总额作分母，比例约为 $5.77\%$，结论不变。
  \item 以表中资本和金融账户差额作为净资本流入 NCI，有 $\text{NCI}=-20,000,000$。理论上
  \[
  NX=-\text{NCI}=20,000,000.
  \]
  \item 第一种名义汇率表示“1 单位外币可兑换多少单位本币”，因此从 X 国视角有 $E=2$ X 币/美元。第一种真实汇率可写为
  \[
  RER=\frac{E P^*}{P}=\frac{2\times 390}{690}=\frac{780}{690}\approx 1.130.
  \]
  含义是：一件美国饰品折成 X 币后的价格，约等于 $1.13$ 件 X 国饰品的价格。
  \item 新的第一种名义汇率为 $E'=3$ X 币/美元，新的真实汇率为
  \[
  RER'=\frac{E'P^{*'}}{P'}=\frac{3\times 310}{790}=\frac{930}{790}\approx 1.177.
  \]
  真实汇率从约 $1.130$ 上升到约 $1.177$。按课堂的第一种真实汇率口径，真实汇率上升表示本国产品相对外国产品更便宜、国际竞争力增强。因此 X 国该产品的国际竞争力上升；名义贬值带来的竞争力提升超过了本国通胀和美国价格下降带来的抵消效应。
\end{enumerate}
\ansmath{\text{净误差与遗漏}=-30,000,000,\quad RER\approx 1.130,\quad RER'\approx 1.177}
\end{solution}
\end{question}

\textbf{七、可贷资金市场与净外汇市场（7 分）}

\begin{question}
在一个持续存在显著贸易顺差的国家（Y 国），政府在考虑推行投资税收减免政策，以补贴国内投资。请回答如下问题：
\begin{enumerate}[label=(\alph*),leftmargin=2.5em]
  \item 该政策将会如何影响国内储蓄、国内投资、净资本流出（NCO）、真实利率、真实汇率和国际贸易差额？请使用下图的格式，结合图像模型，进行解释。（5 分）
\end{enumerate}

\begin{center}
\includegraphics[width=0.82\linewidth]{\finalsamplefigdir/经济学原理（II）-期末样题-fig03.png}

\textbf{图 3：可贷资金市场--净资本流出--净外汇市场三面板图模版}
\end{center}

\begin{enumerate}[label=(\alph*),leftmargin=2.5em,resume]
  \item Y 国许多大型出口商的代表公开反对该政策。这是为什么？请基于上述回答进行解释。（2 分）
\end{enumerate}
\begin{solution}
\begin{enumerate}[label=(\alph*),leftmargin=2.5em]
  \item 投资税收减免提高任一真实利率下的投资意愿，在可贷资金市场中表现为需求曲线右移。均衡真实利率上升，沿着储蓄曲线国内储蓄量增加；由于投资需求曲线本身右移，国内投资也增加。真实利率上升使本国资产更有吸引力，净资本流出 NCO 下降。在净外汇市场中，Y 国有贸易顺差，NCO 是 Y 币净供给来源；NCO 下降使 Y 币供给曲线左移。若按图中“第二种真实汇率”作纵轴，均衡真实汇率上升，代表 Y 币真实升值；等价地，按第一种真实汇率口径则是下降。本币真实升值会抑制出口、刺激进口，所以净出口和贸易顺差下降。
  \item 出口商反对该政策，是因为该政策通过提高真实利率、减少 NCO 并推动本币真实升值，使本国产品在国际市场上变贵。出口竞争力下降，出口量和利润可能减少；对于原本依赖贸易顺差的出口部门尤其不利。
\end{enumerate}
\end{solution}
\end{question}

\textbf{八、总需求与总供给模型（7 分）}

\begin{question}
请根据图 4，回答下列问题。假设总需求曲线目前的位置是 $AD_1$，短期总供给曲线目前的位置是 $SRAS_1$。

\begin{center}
\includegraphics[width=0.42\linewidth]{\finalsamplefigdir/经济学原理（II）-期末样题-fig04.png}

\textbf{图 4：总需求与总供给}
\end{center}

\begin{enumerate}[label=(\alph*),leftmargin=2.5em]
  \item 目前经济的产出缺口是什么类型的缺口？央行什么样的公开市场操作可以“关上”这个缺口？（2 分）
  \item 请举一个“自动稳定器”的例子，说明该“稳定器”通过总需求发挥作用的机制。（2 分）
  \item 假如宏观经济目前处于 $AD_2$ 和 $SRAS_1$ 的交点，且央行打算降低利率，短期内产出和物价水平会有什么变化？长期呢？用长期的结果讨论“货币中性”的含义。（3 分）
\end{enumerate}
\begin{solution}
\begin{enumerate}[label=(\alph*),leftmargin=2.5em]
  \item 当前在 $AD_1$ 与 $SRAS_1$ 的交点，产出为 $Y_1<Y_2=Y_{fe}$，因此是衰退性缺口或负产出缺口。央行应进行扩张性公开市场操作，即买入债券，增加货币供给、降低利率，推动总需求右移以关闭缺口。
  \item 例子：失业保险。经济下行时失业人数上升，失业保险支出自动增加，居民可支配收入下降幅度被缓冲，消费减少得更少，从而稳定总需求。累进所得税也类似：收入下降时平均税负自动下降，缓冲可支配收入和消费。
  \item 若初始在 $AD_2$ 与 $SRAS_1$ 的交点，经济大致处于长期均衡。央行降息会使总需求右移到类似 $AD_3$ 的位置。短期内，经济沿 $SRAS_1$ 移动，产出高于潜在产出，物价水平上升。长期中，工资和其他投入价格上升，短期总供给左移至类似 $SRAS_2$，产出回到潜在产出 $Y_{fe}$，但物价水平永久更高。长期结果说明货币中性：货币政策在长期主要影响价格水平等名义变量，不改变实际 GDP、自然失业率等真实变量。
\end{enumerate}
\end{solution}
\end{question}

\textbf{九、凯恩斯交叉与 AD--AS 的图像模型（7 分）}

\begin{question}
假设宏观经济正处在长期均衡点，实际 GDP 与潜在 GDP 相等。现在，由于人们对未来 GDP 的预期降低，计划投资减少。请回答如下问题：
\begin{enumerate}[label=(\alph*),leftmargin=2.5em]
  \item 请结合 AD--AS 的图像模型阐释新的短期均衡。（2 分）
  \item 请用图像模型展示总支出曲线的变化，并解释为什么会发生这些变化。（3 分）
  \item 请用图像模型展示宏观经济在长期会如何调整到新的长期均衡，并解释其调整的原因。（2 分）
\end{enumerate}
\begin{solution}
\begin{enumerate}[label=(\alph*),leftmargin=2.5em]
  \item 预期未来 GDP 降低会使计划投资减少，进而使总需求曲线左移。短期总供给曲线不立即调整，新的短期均衡出现在原长期均衡左侧，实际产出低于潜在产出，物价水平下降，经济出现衰退性缺口。
  \item 在凯恩斯交叉图中，计划投资减少属于自主支出的下降，因此总支出曲线 $AE=C+I+G+NX$ 整体向下移动。由于均衡条件是 $Y=AE$，新的交点位于更低的实际 GDP。收入下降还会通过消费函数进一步压低消费，产生乘数效应，所以总产出下降幅度通常大于最初的投资下降幅度。
  \item 长期中，衰退性缺口带来失业和闲置产能，名义工资和投入价格下行，短期总供给曲线逐渐右移。经济最终回到 LRAS 上，实际 GDP 回到潜在 GDP，但物价水平低于最初水平。也可以从凯恩斯交叉理解：较低物价提高实际货币余额并降低利率，刺激投资和总支出，直到产出恢复到潜在水平。
\end{enumerate}
\end{solution}
\end{question}

\textbf{十、AD--AS 的数学模型（8 分）}

\begin{question}
假设 Z 国经济学家搜集或估计了该国 2025 年的如下信息：
\begin{center}
\begin{tabular}{rl}
法定准备金率 & $RRR=0.07$\\
货币需求函数 & $M^d=20{,}000-40{,}000r$\\
投资函数 & $I=400-1{,}000r$\\
储蓄函数 & $S=300+1{,}000r$\\
总花费 & $AE=C+G+I+NX$\\
消费函数 & $C=2{,}400+0.4(Y-T+TR)-100P$\\
政府支出 & $G=300$\\
进口 & $M=1{,}000+0.2(Y-T+TR)$\\
出口 & $X=1{,}800$\\
税收 & $T=0.1Y$\\
支付转移 & $TR=0.05Y$\\
总需求 & $AD=AE=Y=C+G+I+NX$\\
长期总供给 & $LRAS=Y^{fe}=3{,}000$\\
短期总供给 & $SRAS: Y=500P-1{,}500$\\
\end{tabular}
\end{center}
请根据如上信息，回答如下问题。
\begin{enumerate}[label=(\alph*),leftmargin=2.5em]
  \item 若可贷资金市场均衡、货币市场出清，请写出 AD 的表达式。（3 分）
  \item 在短期，均衡的真实 GDP 和物价水平（$P$）是多少？请展示计算过程。（2 分）
  \item 政府应该给经济“降温”还是“打气”？如果政府用财政政策，应该如何改变政府支出 $G$ 以实现失业率回归“自然”？（3 分）
\end{enumerate}
\begin{solution}
先由可贷资金市场均衡得到真实利率：
\[
S=I\Rightarrow 300+1000r=400-1000r\Rightarrow r=0.05.
\]
因此
\[
I=400-1000\times 0.05=350.
\]
题目只给出 $RRR$ 而没有给出基础货币，所以 $RRR$ 本身不能单独决定货币供给；在本题中，货币市场出清可理解为央行和银行体系提供与 $r=0.05$ 相容的货币供给，即 $M^s=M^d=18,000$。若额外假设没有现金漏损和超额准备金，则相容的准备金或基础货币为 $18,000\times 0.07=1,260$，但这一数值不改变 AD 的表达式。

又有
\[
Y-T+TR=Y-0.1Y+0.05Y=0.95Y,
\]
所以
\[
C=2400+0.4\times 0.95Y-100P=2400+0.38Y-100P,
\]
\[
NX=X-M=1800-\bigl(1000+0.2\times 0.95Y\bigr)=800-0.19Y.
\]
因此总需求满足
\[
Y=C+G+I+NX=(2400+0.38Y-100P)+300+350+(800-0.19Y),
\]
\[
0.81Y=3850-100P,
\]
即
\[
AD:\quad Y=\frac{3850-100P}{0.81}\approx 4753.09-123.46P.
\]

短期均衡由 $AD=SRAS$ 决定：
\[
\frac{3850-100P}{0.81}=500P-1500.
\]
两边乘以 $0.81$：
\[
3850-100P=405P-1215,
\]
\[
5065=505P\Rightarrow P\approx 10.03.
\]
于是
\[
Y=500\times 10.03-1500\approx 3514.85.
\]
因为 $Y\approx 3514.85>Y^{fe}=3000$，经济处在高于潜在产出的过热状态，政府应该给经济“降温”。

若只通过改变 $G$ 来使短期均衡回到自然水平，则在目标点 $Y=3000$ 时，由 SRAS 得
\[
3000=500P-1500\Rightarrow P=9.
\]
令新的政府支出为 $G'$，AD 变为
\[
Y=(2400+0.38Y-100P)+G'+350+(800-0.19Y),
\]
即
\[
0.81Y=3550+G'-100P.
\]
代入 $Y=3000$、$P=9$：
\[
0.81\times 3000=3550+G'-900,
\]
\[
2430=2650+G'\Rightarrow G'=-220.
\]
因此若机械地只调整政府支出，需要
\[
\Delta G=G'-G=-220-300=-520.
\]
这个数值意味着仅靠 $G$ 调整需要大幅紧缩，甚至把政府支出降为负值，在现实中不可行；政策含义是应采取强紧缩财政，或配合增税、减少转移支付、紧缩货币政策等方式共同“降温”。
\ansmath{AD:\;Y\approx 4753.09-123.46P,\quad P\approx 10.03,\quad Y\approx 3514.85,\quad \Delta G=-520}
\end{solution}
\end{question}

